% Source: 03 - Wed Sep 30 2026.pdf, page 9. Content remains in source order.
\sourcepage{03}{9}

But so far we only created an algorithm. We need to prove that the algorithm actually decides $L_1$ correctly. We need to prove that $L_1=L_{A_{L_1}}$.
We can do it by proving both that  ($L_1 \subseteq L_{A_{L_1}}$ and $L_{A_{L_1}} \subseteq L_1$):


\[
\begin{aligned}
x\in L_1&\Rightarrow y=f(x)\in L_2\Rightarrow A_{L_1}(x)=A_{L_2}(y)=1\\
&\Rightarrow x\in L_{A_{L_1}}\qquad L_1\subseteq L_{A_{L_1}}.
\end{aligned}
\]
\[
\begin{aligned}
x\in L_{A_{L_1}}&\Rightarrow A_{L_1}(x)=A_{L_2}(y)=A_{L_2}(f(x))=1\\
&\Rightarrow f(x)\in L_2\Rightarrow x\in L_1\qquad L_{A_{L_1}}\subseteq L_1.
\end{aligned}
\]
We proved $L_1\subseteq L_{A_{L_1}}$ and $L_1\supseteq L_{A_{L_1}}\Rightarrow L_1=L_{A_{L_1}}$.

$A_{L_1}$ is a poly-time algorithm deciding $L_1$, therefore $L_1\in\classP$.

\textbf{REMARK:} The proof uses both implications $x\in L_1\Longleftrightarrow f(x)\in L_2$.

Poly-time reducibility $\polyred$ is a relation between languages.

\textbf{Homework! Prove that:}
\[
\underbrace{L\polyred L}_{\text{reflexivity}},\qquad
\underbrace{(L_1\polyred L_2)\land(L_2\polyred L_3)\Rightarrow(L_1\polyred L_3)}_{\text{transitivity}}.
\]
